% !TeX root = ../main.tex

\chapter{Marco Teórico}

\section{Misión Gaia}
    La misión Gaia, desarrollada por la Agencia Espacial Europea (ESA), es un 
    observatorio espacial de astrometría lanzado en diciembre de 2013 con el 
    objetivo de cartografiar la Vía Láctea en tres dimensiones con una precisión 
    sin precedentes. Su propósito fundamental es la construcción de un censo 
    estereoscópico y cinemático de aproximadamente mil millones de estrellas, 
    lo que corresponde a cerca del 1\% de la población estelar galáctica 
    \cite{Perryman2002}. Para ello, Gaia mide con muy alta precisión las 
    posiciones ecuatoriales ($\alpha$, $\delta$), los paralajes ($\bar{\omega}$) 
    y los movimientos propios ($\mu_\alpha$, $\mu_\delta$) de cada fuente, 
    complementados con información fotométrica multibanda y, para un subconjunto 
    de estrellas brillantes, velocidades radiales.

    En comparación con su predecesora, la misión Hipparcos que catalogó del 
    orden de $10^5$ estrellas con una precisión significativamente menor, Gaia 
    extiende su alcance a casi dos mil millones de fuentes, incrementando la 
    exactitud astrométrica hasta en dos órdenes de magnitud \cite{ESA2013facts}. 
    Este salto cualitativo permite no solo reconstruir la estructura 
    tridimensional de la Galaxia, sino también abordar problemas fundamentales 
    relacionados con su formación, evolución dinámica y la distribución de la 
    materia oscura.

    Gaia ha aportado uno de los conjuntos de datos astronómicos más completos y 
    homogéneos disponibles en la actualidad. En la segunda publicación de datos 
    (\textit{Gaia Data Release} 2, DR2), se incluyeron aproximadamente 1.7 mil 
    millones de fuentes con fotometría en la banda $G$, de las cuales alrededor 
    de 1.3 mil millones cuentan con medidas de paralaje y movimiento propio 
    \cite{Brown2018}. Adicionalmente, se proporcionó fotometría en las bandas 
    $G_{BP}$ y $G_{RP}$ para cerca de 1.4 mil millones de estrellas, así como 
    velocidades radiales para unos 7 millones de objetos brillantes 
    \cite{Brown2018}. Para una fracción significativa de estas fuentes, 
    se derivaron parámetros físicos estelares como temperaturas efectivas, 
    extinción interestelar y radios a partir de espectros de baja resolución.

    Las versiones posteriores de los datos (EDR3 y DR3, publicadas entre 2020 y 
    2022) han refinado sustancialmente la precisión de estas mediciones y han 
    incorporado nuevas clases de objetos, incluyendo estrellas variables, 
    asteroides y fuentes extragalácticas, gracias a un mayor intervalo temporal 
    de observaciones.

    La alta precisión astrométrica y fotométrica de Gaia resulta particularmente 
    valiosa para el estudio de cúmulos estelares. Las estrellas miembros de un 
    cúmulo abierto comparten propiedades cinemáticas similares, como paralajes y 
    movimientos propios coherentes, además de definir secuencias bien 
    delimitadas en el diagrama color-magnitud \cite{CastroGinard2019}. En 
    consecuencia, los cúmulos pueden identificarse como sobredensidades en 
    espacios de parámetros multidimensionales, típicamente definidos por la 
    posición en el cielo, el paralaje y el movimiento propio. Diversos estudios 
    han demostrado la eficacia de algoritmos de agrupamiento no supervisado, 
    como DBSCAN, para la detección automática de cúmulos en los catálogos de 
    Gaia \cite{CastroGinard2019}.

    Por ejemplo, Castro-Ginard et al. (2019) aplicaron técnicas de detección 
    basadas en aprendizaje automático a los datos de Gaia DR2 y descubrieron 
    decenas de cúmulos abiertos previamente desconocidos en la región del 
    anticentro galáctico, incrementando significativamente el censo conocido en 
    esa zona \cite{CastroGinard2019}. De manera similar, Cantat-Gaudin et al. 
    (2018) construyeron un catálogo homogéneo de más de mil cúmulos abiertos 
    confirmados a partir de Gaia DR2 \cite{cantat2018gaia}. Y con el nuevo 
    lanzamiento de Gaia DR3, Hunt y Reffert (2023) mejoraron el censo de cúmulos 
    abiertos en la Vía Láctea reconociendo 7167 de estos objetos 
    \cite{hunt2023improving}. En conjunto, estos trabajos evidencian que Gaia ha 
    revolucionado la detección, confirmación y caracterización de cúmulos 
    estelares, permitiendo estudiar su distribución espacial, edades, cinemática 
    interna y procesos evolutivos con un gran nivel de detalle 
    \cite{CastroGinard2019}.

    La determinación de distancias estelares en Gaia se basa fundamentalmente en 
    la medición del paralaje trigonométrico. En su forma más simple, la 
    distancia $d$ (parsecs) se relaciona con el paralaje $\bar{\omega}$ (arcseg) 
    mediante la expresión $d \approx 1 / \bar{\omega}$. No obstante, debido a 
    las incertidumbres observacionales, esta aproximación resulta inadecuada 
    cuando los errores relativos del paralaje son significativos o cuando se 
    obtienen valores de paralaje pequeños o incluso negativos.

    En este contexto, Luri et al. (2018) subrayan la necesidad de emplear 
    enfoques estadísticos, particularmente métodos bayesianos, para inferir 
    distancias fiables a partir de los datos de Gaia \cite{Luri2018}. Estos 
    métodos combinan la información observacional con distribuciones previas 
    basadas en modelos astrofísicos de la distribución estelar en la Galaxia, 
    permitiendo explotar de manera óptima incluso mediciones con grandes 
    incertidumbres. Un ejemplo destacado es el catálogo de distancias bayesianas 
    publicado por Bailer-Jones et al. (2018) \cite{bailer2018estimating}, 
    ampliamente utilizado en estudios galácticos y extragalácticos.

    La precisión alcanzada por Gaia es excepcional. Para estrellas brillantes 
    ($G < 15$), los errores típicos en posición son del orden de decenas de 
    microsegundos de arco \newline \cite{ESA2013facts}, lo que se traduce en 
    errores relativos en distancia extremadamente pequeños para objetos 
    cercanos. Incluso para estrellas situadas a distancias de varios 
    kiloparsecs, las incertidumbres en paralaje permiten estimaciones de 
    distancia con precisiones del orden del 10-20 \% \cite{ESA2013facts}. Estas 
    cifras superan ampliamente las capacidades de misiones anteriores, como 
    Hipparcos, cuyas precisiones se limitaban al nivel de milisegundos de arco. 
    Gracias a esta exactitud, Gaia proporciona distancias muy precisas para 
    numerosos cúmulos estelares galácticos, obtenidas a partir del paralaje 
    medio de sus miembros. Esto permite restringir de forma robusta parámetros 
    fundamentales como la edad, la masa y la estructura interna de los cúmulos.

\section{Cúmulos Estelares}
    Estos son grupos gravitacionalmente ligados de decenas a miles de estrellas 
    nacidas simultáneamente de la misma nube molecular. Existen cúmulos abiertos 
    que se caracterizan por ser jóvenes de baja masa y se encuentran en el disco 
    galáctico, y cúmulos globulares que son cúmulos antiguos, masivos y se 
    encuentran en el halo. Su estudio aborda la identificación de miembros, 
    propiedades físicas comunes y la dinámica interna. 

    Para determinar qué estrellas pertenecen al cúmulo y cuáles al campo de la 
    Galaxia clásicamente se usan métodos paramétricos: se ajustan distribuciones 
    de velocidad y posición p. ej. mezcla de gaussianas para asignar 
    probabilidades de membresía (el método de Sanders 1971 
    \cite{sanders1971improved} es un ejemplo histórico). Con los datos modernos 
    de Gaia se emplean cada vez más algoritmos no paramétricos y de aprendizaje 
    automático. Un método ampliamente usado es UPMASK, que combina clustering 
    por k-medias con pruebas de significancia en el espacio multivariado 
    (posición, paralaje, movimiento propio) \cite{Elsanhoury2026}. Por ejemplo, 
    Cantat-Gaudin et al. (2018) aplicaron UPMASK a los datos de Gaia DR2 para 
    miles de cúmulos conocidos \cite{cantat2018gaia}. UPMASK demuestra alta 
    puridad y completitud al separar estadísticamente las estrellas reales del 
    cúmulo del fondo estelar \cite{Elsanhoury2026}\cite{cantat2018gaia}. Otros 
    enfoques modernos incluyen métodos bayesianos, p. ej. análisis bayesiano de 
    campo estelar, y algoritmos de agrupamiento como DBSCAN. En síntesis, la 
    combinación de astrometría (coordenadas y movimientos propios) y fotometría 
    en el espacio de fase permite identificar con robustez los miembros reales 
    del cúmulo en presencia de contaminación de campo \cite{Elsanhoury2026} 
    \cite{cantat2018gaia}.

    Por otra parte, los cúmulos estelares comparten propiedades físicas 
    medibles: masa total, radios característicos, densidad estelar, edad y 
    metalicidad. La masa total suele estimarse sumando las masas individuales de 
    estrellas miembros (usando modelos M/L), o bien dinámicamente mediante el 
    teorema del virial. Según este último, un cúmulo en equilibrio obedece $2T + 
    \Omega = 0$, donde $T$ y $\Omega$ son energía cinética y potencial, lo que 
    lleva a una relación aproximada $M \thicksim \sigma^2 R/G$, donde $\sigma$ 
    es la dispersión de velocidad y $R$ es el radio característico 
    \cite{Elsanhoury2026}. Por ejemplo, Elsanhoury et al. (2026) usan 
    dispersiones propias y radiales para derivar masas totales de cúmulos 
    abiertos, encontrando desde unos $5 \times 10^2 \ \text{M}_\odot$ hasta 
    $\thicksim 3 \times 10^3 \ \text{M}_\odot$ en sus sistemas estudiados 
    \cite{Elsanhoury2026}.

    La densidad interna y distribución radial se modelan con perfiles teóricos 
    p. ej. modelos de King o Plummer que ajustan el brillo superficial 
    observado. El radio tidal o de Jacobi es la distancia de equilibrio con el 
    campo galáctico, estimado por equilibrio gravitacional. Para un cúmulo de 
    masa $m$ orbitando a distancia $R$ de una galaxia de masa $M$, la fórmula de 
    Spitzer aproxima $r_t \approx R(m/3M)^{1/3}$ \cite{Renaud2018}. Este radio 
    indica hasta dónde domina la gravedad del cúmulo sobre las mareas 
    galácticas. En la práctica, parámetros derivados de perfiles de brillo como 
    los radios central, de medio perfil y tidal se extraen mediante ajuste de 
    modelos isotrópicos de equilibrio.

    El diagrama Hertzsprung-Russell (diagrama color-magnitud) es herramienta 
    clave. Graficando magnitud absoluta versus color o temperatura, el cúmulo 
    revela una secuencia principal común para todas sus estrellas integrantes, 
    ya que estas nacieron casi a la vez, así como ramas evolutivas de gigantes 
    jóvenes o viejas. Ajustando isócronas teóricas se determina la edad y en 
    ocasiones metalicidad del cúmulo. El ajuste puede realizarse mediante 
    técnicas bayesianas o MCMC tal como en Elsanhoury et al. (2026), donde 
    usando Gaia DR3 y UPMASK para limpieza, obtuvieron edades entre $\log \tau = 
    7.0 - 8.15$ para cúmulos jóvenes \cite{Elsanhoury2026}. Por otra parte, 
    comprendiendo qué parte de los diagramas color-color de un cúmulo 
    corresponde a su respectivo diagrama color-magnitud, se facilita la 
    caracterización de la composición estelar de los cúmulos ya que los 
    diagramas color-color no dependen de la distancia como sí lo hacen los 
    diagramas color-magnitud.

    Los cúmulos se forman en nubes moleculares gigantes. Los cúmulos embebidos 
    (recién nacidos dentro de gas y polvo) son abundantes, sin embargo, tras 
    algunas decenas de Myr muchos se disgregan (el gas es expulsado por 
    retroalimentación masiva). Estudios como Lada \& Lada (2003) muestran que la 
    mayoría de las estrellas nacen en agrupamientos embebidos, pero pocos 
    sobreviven como cúmulos abiertos maduros \cite{Lada2003}. El efecto de marea 
    como supernovas o viento estelar y la rotación diferencial galáctica 
    determinan la vida útil. Un cúmulo inicialmente subvirial puede ``rebotar'' 
    al expulsar gas; solo con una alta eficiencia de formación ($\gtrsim 50 \%$) 
    se mantiene ligado. De esta forma, la evolución temprana, hasta $\thicksim 
    100 \ \text{Myr}$, suele ser dinámicamente activa: los cúmulos jóvenes a 
    menudo no están aún en equilibrio virial completo. Conforme envejecen, el 
    cúmulo alcanza equipartición de energía y puede sufrir colapso del núcleo, 
    expulsión lenta de estrellas, etc. \cite{Spitzer1987}\cite{heggie2003book}.

    La función de luminosidad de un cúmulo describe el número de estrellas en 
    función de su brillo, típicamente se mide combinando fotometría profunda con 
    corrección de miembros. Esta está directamente relacionada con la función de 
    masa inicial (IMF). Históricamente, Salpeter (1955) propuso una IMF de ley 
    de potencia $\xi(M) \varpropto M^{-\alpha}$ con $\alpha \approx 2.35$, es 
    decir una pendiente $-2.35$ en la escala $\log N$ vs $\log M$ 
    \cite{salpeter1955luminosity}. En cúmulos abiertos se suele observar una 
    pendiente próxima a Salpeter en masas medias-altas. Por ejemplo, Elsanhoury 
    et al. (2026) hallan pendientes de IMF $\thicksim -2.26$, muy similares a 
    $-2.35$, en cúmulos jóvenes \cite{Elsanhoury2026}. A bajas masas ($< 0.5 \ 
    \text{M}_\odot$) la IMF se aplana (modelo de Kroupa 2001 
    \cite{kroupa2001variation} o Chabrier 2003 \cite{chabrier2003galactic}). 
    Estudios de cúmulos masivos revelan que, a pesar de distintas condiciones, 
    la IMF parece casi universal (Kroupa 2002 \cite{kroupa2002initial}; 
    Bastian et al. 2010 \cite{bastian2010universal}).

    La función de luminosidad resulta de combinar la IMF con la función de 
    masa-luminosidad dada por la evolución estelar y la edad. Es conveniente en 
    cúmulos jóvenes, es decir, aún con mayoría de estrellas en secuencia 
    principal, donde la transformación masa-brillo es monótona. Con estas 
    funciones se infieren la masa total integrada y la población estelar. En 
    trabajos recientes, se modela simultáneamente la IMF y la función de 
    luminosidad usando datos fotométricos profundos y algoritmos bayesianos, 
    para descartar estrellas de campo y estimar la pendiente $\alpha$ y la masa 
    total \cite{Elsanhoury2026}\cite{cantat2018gaia}.

    Ahora, para las condiciones dinámicas internas y equilibrio, Desde el punto 
    de vista dinámico, un cúmulo puede aproximarse como un sistema 
    auto-gravitante que busca un estado de equilibrio dinámico. Para cúmulos 
    antiguos, p. ej. globulares, se asume a menudo un equilibrio virial, donde 
    la energía cinética neta es la mitad en magnitud de la energía potencial 
    (Teorema del virial). En dicho caso el cúmulo es isotrópico y estacionario. 
    En cambio, cúmulos jóvenes o que estén disgregándose pueden estar fuera de 
    equilibrio; su virialización depende de factores como la expulsión de gas, 
    colisiones y perturbaciones externas.

    El tiempo de relajación $t_\text{relax}$ es el lapso en que las 
    interacciones de dos cuerpos (colisiones gravitacionales débiles) 
    redistribuyen la energía cinética interna. Aproximadamente $t_\text{relax} 
    \thicksim (0.1N /\ln N)t_\text{cross}$, con $N$ el número de estrellas y 
    $t_\text{cross}$ el tiempo de cruce. Si la edad del cúmulo supera 
    $t_\text{relax}$, alcanza equipartición parcial y se dice que está relajado. 
    En Elsanhoury et al. (2026) se encuentra que sólo los cúmulos más viejos del 
    estudio, NGC 2301 y NGC 2384b habían relajado su población; los más jóvenes 
    aún evolucionan dinámicamente \cite{Elsanhoury2026}. Esto concuerda con la 
    idea de que la mayoría de los cúmulos abiertos persisten pocos cientos de 
    Myr antes de evaporarse hacia el campo galáctico.

    Por otra parte, según la densidad, los cúmulos pueden ser colisionales 
    (frecuentes interacciones de trayectorias estelares) o no colisionales. Los 
    cúmulos globulares densos tienen $t_\text{relax}$ más corto que la edad del 
    universo, por lo que la dinámica colisional es esencial (core-collapse, 
    procesos de Spitzer, etc.). Muchos cúmulos abiertos, al ser poco masivos, 
    son casi no colisionales en su halo. La frontera entre ambos regímenes se 
    ilustra comparando $t_\text{relax}$ con la edad: si $t_\text{relax} \ll$ 
    edad, el cúmulo tiende al equilibrio térmico, es colisional; si 
    $t_\text{relax} \gg $ edad, sigue más un flujo libre colisional, es no 
    colisional. Esta distinción determina la validez de la aproximación de 
    Boltzmann no colisional: en cúmulos con larga relajación, la ecuación de 
    Boltzmann sin término de colisión (Jeans) describe bien el equilibrio medio; 
    si no, hay que añadir términos de difusión de Boltzmann o simulaciones 
    N-cuerpos completas (Heggie \& Hut 2003) \cite{heggie2003book}.

    Los cúmulos se caracterizan por radios típicos:

    \begin{itemize}
        \item \textbf{Radio central o core} ($r_c$): donde la densidad cae a la 
            mitad del centro.

        \item \textbf{Radio de mitad de masa} ($r_h$): radio que encierra la 
            mitad de la masa total.

        \item \textbf{Radio tidal} ($r_t$): límite externo impuesto por las 
            mareas galácticas; en modelos de King (1962) aparece como radio 
            donde la densidad proyectada alcanza cero. Matemáticamente, la 
            aproximación clásica para órbita circular en potencial esférico da 
            $r_t \simeq R(m/3M)^{1/3}$ \cite{Spitzer1987}. Este radio tidal 
            también llamado radio de Jacobi, define el volumen de influencia 
            gravitacional del cúmulo. En la práctica, se ajustan modelos 
            teóricos como King, Plummer, o modelos modernos ``limepy'' 
            \cite{Cheng2023}, a perfiles de brillo superficial para inferir 
            estos radios y la concentración (cociente $r_t/r_c$). 
    \end{itemize}

    En las observaciones, se utilizan perfiles de densidad o brillo radial: los 
    modelos isotrópicos clásicos (King 1966) ajustan bien los cúmulos 
    globulares, mostrando un núcleo plano y caída acentuada cerca de $r_t$. 
    El modelo de Plummer (1911) es otro perfil simétrico con densidad $\rho(r) 
    \varpropto \left(1 + (r/a)^2\right)^{-5/2}$, donde $a$ es el radio 
    característico, que no se trunca naturalmente, aunque puede aproximar 
    regiones internas. Según Cheng \& Jiang (2023), su familia de modelos 
    ``limepy'' incorpora King y Plummer como casos especiales: el modelo de 
    King surge restando una energía de corte al DF isotrópico \cite{Cheng2023}, 
    y el Plummer corresponde a un caso particular $g=3.5$ con masa finita e 
    extensión infinita \cite{Cheng2023}. Estas modelizaciones dinámicas permiten 
    relacionar observaciones proyectadas con distribuciones de masa 
    tridimensional.

    La masa de un cúmulo se estima cinemáticamente midiendo la dispersión de 
    velocidades de las estrellas, tanto radiales como tangenciales, y aplicando 
    los esquemas anteriores (teorema del virial o ecuaciones de Jeans). En la 
    forma más simple, es decir, asumiendo un cúmulo esférico o isotrópico, el 
    teorema del virial da $M \approx \frac{5\sigma^2 r_h}{G}$ o similar, donde 
    el factor de 5 depende del perfil. Alternativamente, integrando la ecuación 
    de Jeans radial, Ecuación~\eqref{eq:jeans_esferica}, se obtiene la masa 
    interior a $r$. Al conocer $\sigma(r)$ y $\rho(r)$, se inferiría el 
    potencial $\Phi(r)$ y por ende la masa. Sin embargo, sin $\beta$ se suele 
    asumir isotropía ($\beta = 0$), Ecuación~\eqref{eq:jeans_integral}. En la 
    práctica se comparan perfiles de velocidad con modelos como King, Plummer u 
    otros de distribución, que que fijan masa total y M/L. Por ejemplo, 
    Cantat-Gaudin et al. usaron UPMASK más isocronas para encontrar masa 
    integrada y distribución de masa \cite{cantat2018gaia}, y Elsanhoury et al. 
    obtuvieron masas totales directamente de tales ajustes MCMC de Gaia 
    \cite{Elsanhoury2026}. Los cúmulos más masivos o densos requieren explicar 
    por qué su valor M/L dinámico puede exceder el de la población de estrellas 
    visibles, esto puede ser por la presencia de enanas blancas, estrellas 
    binarias, etc.

    Una técnica actual es ajustar conjuntamente brillo superficial y perfil de 
    dispersión de velocidad con modelos de DF (King/Michie o limepy). Cheng \& 
    Jiang (2023) aplicaron precisamente esta estrategia con limepy a 18 cúmulos 
    globulares, determinando masa total, radio de mitad de masa y relaciones M/L 
    que describen bien los datos \cite{Cheng2023}. La modelización detallada 
    permite inferir también la anisotropía radial que es el exceso de órbitas 
    radiales frente a tangenciales, parámetro importante para la estabilidad 
    futura del cúmulo.

    Para describir la estructura interna de cúmulos en equilibrio se usan 
    modelos de distribución de velocidades. Los modelos de King (1962) suponen 
    una distribución de Maxwell con energía truncada: la función de distribución 
    dependiente de la energía específica $E$ toma la forma $f(E) \varpropto 
    \exp{(-E/\sigma^2)}$ para $E < 0$ y $f = 0$ para $E > 0$, donde la energía 
    cero corresponde a la energía en el radio tidal. Esto genera un perfil de 
    densidad esférico con núcleo plano y corte abrupto en $r_t$. Como explican 
    Cheng \& Jiang (2023), el modelo de King se obtiene al restar sucesivas 
    constantes del DF para truncar la distribución hasta energía cero 
    \cite{Cheng2023}. En cambio, el modelo de Plummer (1911) es un caso de 
    esfera isotrópica sin truncación; su densidad $\rho_P(r) = \frac{3M}{4\pi 
    a^3}\left(1 + r^2/a^2\right)^{-5/2}$ se extiende infinitamente sin llegar a 
    cero \cite{Cheng2023}. Cheng \& Jiang señalan que la familia “limepy” 
    generaliza ambos: el caso de Plummer (finita la masa pero sin corte) 
    equivale a un parámetro de truncamiento $g = 3.5$ \cite{Cheng2023}, mientras 
    que $g = 1$ recupera el perfil de King isotrópico \cite{Cheng2023}.

    Estas soluciones estacionarias surgen de imponer la ecuación de Jeans en 
    equilibrio. En su forma radial simplificada, la ecuación de Jeans liga el 
    gradiente de la presión cinética ($\rho\sigma_r^2$) con la pendiente del 
    potencial gravitarorio. Así, dado un perfil de densidad $\rho(r)$, la 
    dispersión $\sigma(r)$ viene determinada por la masa interior. Los modelos 
    mencionados son casos particulares de soluciones auto-consistentes de dicha 
    ecuación bajo supuestos de isotropía o anisotropía. En la práctica se 
    utiliza software de ajuste que, dado perfiles observados de brillo y 
    dispersión, explora parámetros como la concentración, radio de anisotropía y 
    masa total \cite{Cheng2023}.

\section{El modelo de King}
    \subsection{Contexto Histórico y Motivación del Modelo}
        El estudio dinámico de los cúmulos globulares ha sido siempre un área de 
        gran interés en la astrofísica debido a la riqueza estelar y simetría de 
        estos sistemas. Históricamente, uno de los primeros intentos por 
        describir su estructura matemática fue la ley de Jeans, la cual proponía 
        teóricamente que las regiones externas de un cúmulo debían presentar una 
        densidad estelar decreciente proporcional a $r^{-4}$ (lo que se traduce 
        en $r^{-3}$ al proyectarse en el plano del cielo). Sin embargo, el 
        trabajo de Ivan King en 1962 demostró que la ley de Jeans partía de 
        asunciones teóricas incorrectas, como suponer una distribución de 
        velocidades isotrópica en todo punto del cúmulo y emplear una expansión 
        matemática injustificada para el potencial gravitacional.

        Gracias al uso de observaciones fotográficas profundas, como las 
        obtenidas con la cámara Schmidt del Observatorio de Palomar, se hizo 
        evidente que la densidad estelar no decae asintóticamente hacia el 
        infinito, sino que cae abruptamente a cero en un límite físico finito. 
        Este comportamiento empírico motivó la formulación del modelo de King, 
        el cual unifica la descripción del denso núcleo del cúmulo y su frontera 
        truncada por efectos gravitacionales externos \cite{king1962structure}.

    \subsection{Fundamentos Físicos y Parámetros del Modelo}
        El modelo de King se fundamenta en la idea de que la estructura de un 
        cúmulo estelar no requiere de matemáticas arbitrarias, sino que está 
        regida por las mínimas condiciones impuestas por su entorno físico e 
        historial dinámico. Este modelo asume que todos los cúmulos alcanzan un 
        cuasi-equilibrio derivado de procesos de relajación (mezcla inicial y 
        encuentros estelares) que producen efectos casi idénticos en su 
        arquitectura.

        Físicamente, el modelo se describe de manera invariante empleando 
        únicamente tres parámetros fundamentales, el número mínimo permitido por 
        las restricciones físicas del sistema \cite{king1962structure}:

        \begin{itemize}
            \item \textbf{Factor de escala $k$:} Es un parámetro asociado al 
            número total de estrellas y está directamente relacionado con la 
            densidad superficial central del cúmulo.
            
            \item \textbf{Radio del núcleo ($r_c$):} Actúa como un factor de 
            escala espacial que describe la región interior del cúmulo, y se 
            encuentra intrínsecamente determinado por la energía interna del 
            sistema estelar.

            \item \textbf{Radio límite o de marea ($r_t$):} Representa la 
            distancia desde el centro del cúmulo a la cual la densidad estelar 
            cae a cero. Físicamente, este radio marca el punto donde las 
            estrellas escapan del cúmulo debido a que la fuerza gravitacional 
            del cúmulo es superada por el campo de fuerzas de marea de la Vía 
            Láctea.
        \end{itemize}

        Además, aunque este no es un parámetro independiente, el parámetro de 
        concentración, se define como la razón $c = r_t/r_c$. Este valor 
        clasifica cuán concentrado es un cúmulo en su centro en relación con su 
        extensión total dictada por las fuerzas de marea.

    \subsection{Interpretación de las Distribuciones de Densidad}
        El aporte fundamental del modelo de King radica en la formulación de 
        ecuaciones que describen de forma empírica y coherente la densidad del 
        cúmulo en diferentes dominios dimensionales.

        \subsubsection{Densidad Proyectada o Superficial}
            En la práctica observacional, la astronomía mide la distribución 
            bidimensional de las estrellas proyectadas sobre el plano del cielo. 
            Para ello, King introdujo una ecuación fundamental que rige la 
            densidad superficial estelar $f$ a una distancia radial $r$ del 
            centro del cúmulo:

            \begin{equation}
                f = k \left[ \frac{1}{\sqrt{1 + (r/r_c)^2}} - \frac{1}{\sqrt{1 + 
                (r_t/r_c)^2}} \right]^2
                \label{eq:den_2d}
            \end{equation}

            Esta ecuación integra armónicamente dos comportamientos: en las 
            proximidades del centro ($r \ll r_t$), reproduce una distribución 
            dominada por el radio del núcleo $r_c$; mientras que en las 
            fronteras externas, fuerza a la densidad a caer a cero exactamente 
            cuando $r = r_t$.

        \subsubsection{Densidad Espacial o Tridimensional}
            Para comprender la dinámica real del cúmulo, es imperativo 
            transformar la densidad proyectada bidimensional en una densidad 
            espacial volumétrica $\varphi(r)$. Esta cantidad, que representa la 
            concentración tridimensional de estrellas, se obtiene integrando la 
            densidad superficial. El modelo de King define esta densidad como:

            \begin{equation}
                \varphi(r) = \frac{k}{\pi r_c [1 + (r_t/r_c)^2]^{3/2}} 
                \frac{1}{w^2} \left[ \frac{1}{w} \cos^{-1} w - (1 - w^2)^{1/2} 
                \right]
                \label{eq:den_3d}
            \end{equation}

            Para simplificar la notación geométrica y relacionar la posición $r$ 
            con los radios característicos del cúmulo, se introduce la variable 
            auxiliar adimensional $w$, expresada de la siguiente manera:

            $$w = \sqrt{ \frac{1 + (r/r_c)^2}{1 + (r_t/r_c)^2} }$$

            Al observar la composición de $w$, resulta evidente que en el límite 
            físico del cúmulo, cuando $r$ tiende a $r_t$, el valor de $w$ tiende 
            a $1$. Al evaluar la ecuación de la densidad espacial con $w=1$, el 
            término entre corchetes se anula, garantizando matemáticamente que 
            la densidad volumétrica $\varphi(r)$ sea exactamente cero en la 
            frontera de marea del sistema.

        \subsubsection{Densidad de Franja}
            Adicional a la densidad volumétrica y superficial, King explora la 
            densidad de franja, la cual evalúa el número total de estrellas 
            contenidas a lo largo de una franja transversal de anchura unitaria 
            que cruza completamente el cúmulo a una distancia $r$ del centro. 
            Aunque es una medida menos común hoy en día, proporciona una 
            integral alternativa para comparar perfiles de densidad en recuentos 
            unidimensionales sobre placas fotográficas.

    \subsection{Importancia y Relevancia en la Astrofísica}
        El modelo empírico propuesto por King demostró que, a pesar de sus 
        complejas historias de formación, los cúmulos estelares y por extensión 
        algunas galaxias enanas elípticas son estructuras tan similares como las 
        leyes físicas lo permiten. La simplicidad de la formulación encapsula 
        una verdad física profunda: la morfología gravitacional de un sistema 
        estelar relajado está constreñida enteramente por el número de 
        partículas, la conservación de la energía y las fuerzas de marea 
        externas \cite{king1962structure}.

        Medio siglo después de su publicación, las formulaciones de King 
        continúan siendo el estándar analítico primordial para el ajuste de 
        perfiles de brillo superficial y dispersión de velocidades. El éxito del 
        modelo radica en su capacidad para ofrecer un puente analítico directo 
        entre el recuento observacional de estrellas y las propiedades puramente 
        dinámicas de la materia confinada gravitacionalmente, manteniéndose 
        indispensable para los conductos modernos de arqueología galáctica y las 
        simulaciones computacionales de evolución estelar.

\section{Ecuaciones de Jeans y Dinámica de Sistemas Gravitacionales}
        La dinámica estelar constituye una rama cardinal de la astrofísica que 
        se enfoca en el estudio del movimiento colectivo de un gran número de 
        masas puntuales como estrellas, remanentes estelares o partículas de 
        materia oscura, bajo  la influencia de su campo gravitacional 
        autogenerado. A diferencia de la mecánica celeste clásica, que resuelve 
        analítica o numéricamente las órbitas exactas de un número reducido de 
        cuerpos, a dinámica de sistemas macroscópicos como las galaxias, los 
        halos de materia oscura, los cúmulos globulares y los cúmulos abiertos 
        demanda un enfoque fundamentado en la mecánica estadística y la teoría 
        cinética \cite{BinneyTremaine2008}. 

        El estado microscópico de un sistema estelar de $N$ partículas queda 
        completamente especificado en un espacio de fase de $6N$ dimensiones. 
        Sin embargo, para sistemas con un número suficientemente masivo de 
        componentes, resulta físicamente riguroso y matemáticamente tratable 
        describir el sistema a través de una función de distribución continua 
        en el espacio de fase de una sola partícula, de seis dimensiones, esto 
        es, tres coordenadas espaciales y tres coordenadas de velocidad. Esta 
        función, denotada clásicamente como $f(\mathbf{x}, \mathbf{v}, t)$, 
        representa la densidad de probabilidad o el número de estrellas 
        localizadas en un volumen diferencial del espacio de fase $d^3x d^3v$ 
        centrado en la posición $\mathbf{x}$ y con velocidad $\mathbf{v}$ en un 
        instante de tiempo $t$.
        
        De acuerdo con el teorema de Liouville, en ausencia de disipación y 
        fenómenos de creación o aniquilación, la densidad en el espacio de fase 
        alrededor de un punto representativo del sistema se conserva a lo largo 
        de su trayectoria. Las ecuaciones de Jeans nacen precisamente como un 
        corolario hidrodinámico de esta conservación, proveyendo un andamiaje 
        teórico que conecta la cinemática inobservable tridimensional con las 
        proyecciones observables, constituyendo así el núcleo fundamental de 
        cualquier modelado dinámico de cúmulos estelares. Originalmente 
        derivadas por James Clerk Maxwell en el contexto de la teoría cinética 
        de gases, fueron adaptadas e introducidas formalmente a la astronomía 
        por James Jeans en 1919 \cite{Jeans1919}. Hoy en día, son la herramienta 
        primaria para deducir la masa gravitante total y la estructura orbital 
        de sistemas astronómicos.

        \subsection{La Ecuación de Boltzmann No Colisional}
            Para comprender el origen de las ecuaciones de Jeans, es imperativo 
            partir de la Ecuación de Boltzmann No Colisional (CBE, por sus 
            siglas en inglés), frecuentemente referida en la literatura física 
            como la ecuación de Vlasov.

            Si un sistema evoluciona bajo un potencial gravitacional medio 
            $\Phi(\mathbf{x}, t)$, y si los encuentros gravitacionales locales 
            entre pares de estrellas (colisiones dinámicas) son insignificantes 
            frente al campo global, la evolución temporal de la función de 
            distribución $f(\mathbf{x}, \mathbf{v}, t)$ obedece a la ecuación de 
            continuidad en el espacio de fase \cite{mamon2010kinematic}:

            $$\frac{\partial f}{\partial t} + \mathbf{v} \cdot \nabla_x f - 
            \nabla_x \Phi \cdot \nabla_v f = 0$$

            donde $\nabla_x$ y $\nabla_v$ denotan los operadores gradiente 
            respecto a las coordenadas espaciales y de velocidad, 
            respectivamente. El término $-\nabla_x \Phi$ representa la 
            aceleración gravitacional experimentada por una estrella en la 
            posición $\mathbf{x}$.

            La CBE dictamina que el flujo de probabilidad en el espacio de fase 
            es incompresible. Sin embargo, encontrar soluciones analíticas 
            directas a la CBE, $f(\mathbf{x}, \mathbf{v}, t)$, es un gran 
            desafío, ya que su dominio natural es el espacio hexadimensional 
            \cite{ciotti_2021_jeans}. Es virtualmente imposible medir 
            observacionalmente las seis coordenadas de fase para cada estrella 
            en un sistema denso. Usualmente, las observaciones astronómicas 
            están limitadas a proveer información proyectada bidimensional 
            sobre el plano del cielo y una única componente de velocidad radial 
            o a lo largo de la línea de visión inferida por efecto Doppler 
            espectroscópico. Frente a esta severa restricción,  el tratamiento 
            de momentos integrados sobre el espacio de velocidades, las 
            ecuaciones de Jeans, reduce la dimensionalidad del problema del 
            espacio de fase al espacio de configuración tridimensional regular 
            \cite{courteau2014galaxy}.

        \subsection{Hipótesis Fundamentales y Límites de Aplicabilidad}
            La construcción teórica de las ecuaciones de Jeans en sistemas 
            estelares requiere el cumplimiento de algunas suposiciones 
            fundamentales.  La validez de un modelo de masa derivado para un 
            cúmulo estelar depende enteramente de cuán precisamente la 
            naturaleza del cúmulo se adhiera a las siguientes hipótesis.

            \subsubsection{Sistemas Colisionales, No Colisionales y el Tiempo 
            de Relajación}
                El andamiaje de la CBE y, consecuentemente, de las ecuaciones de 
                Jeans estándar, descansa sobre la presunción de que el sistema 
                es "no colisional". En dinámica gravitacional, una colisión no 
                implica el impacto físico de las superficies estelares, sino una 
                interacción de dos cuerpos lo suficientemente cercana como para 
                que la trayectoria de una estrella se desvíe significativamente 
                respecto a la órbita dictada exclusivamente por el potencial 
                global suavizado \cite{hamilton_2021_kinetic}.

                La escala temporal que demarca la transición entre un régimen 
                gobernado por el potencial medio (no colisional) y uno dominado 
                por interacciones discretas (colisional) es el tiempo de 
                relajación de dos cuerpos, $t_\text{relax}$ . Este se define 
                heurísticamente como el tiempo necesario para que las pequeñas 
                deflexiones acumuladas por encuentros estelares generen un 
                cambio en la velocidad de la estrella del mismo orden de 
                magnitud que su velocidad inicial \cite{hamilton_2021_kinetic}.

                Siguiendo el formalismo clásico de Chandrasekhar 
                \cite{Chandrasekhar1942} y Spitzer \cite{Spitzer1987}, si 
                consideramos una estrella cruzando un sistema de tamaño $R$ con 
                una velocidad típica $v$, el cambio cumulativo en la varianza de 
                su velocidad ortogonal debido a encuentros sucesivos conduce a 
                una caminata aleatoria en el espacio de velocidades. El tiempo 
                de relajación se relaciona íntimamente con el tiempo dinámico o 
                tiempo de cruce, $t_\text{cross} \sim R/v$, y el número total de 
                estrellas del sistema $N$ mediante la relación aproximada:

                $$t_\text{relax} \approx \frac{N}{10 \ln \Lambda} t_\text{cross}
                $$

                donde $\ln \Lambda \approx \ln(N)$ es el logaritmo de Coulomb, 
                un factor proveniente de la integración sobre todos los posibles 
                parámetros de impacto entre un mínimo que produce una deflexión 
                de 90 grados y un máximo típicamente el tamaño del sistema 
                \cite{binney_tremaine_2008}.

                La distinción entre sistemas astronómicos basada en 
                $t_\text{relax}$ es drástica:

                \begin{itemize}
                    \item \textbf{Galaxias y halos de materia oscura:} Con $N 
                    \sim 10^{11}$, los tiempos de relajación exceden los 
                    $10^{14}$ años, siendo varios órdenes de magnitud mayores a 
                    la edad del universo \cite{hamilton_2021_kinetic}. Son 
                    sistemas estrictamente no colisionales donde las ecuaciones 
                    de Jeans operan en su forma más pura.

                    \item \textbf{Cúmulos Globulares:} Compuestos por 
                    $N \sim 10^5 - 10^6$ estrellas confinadas en radios de 
                    escasos parsecs. Su tiempo de relajación oscila entre $10^8$ 
                    y $10^9$ años. Al ser menores que sus edades intrínsecas que 
                    son usualmente $> 10^{10}$ años, os cúmulos globulares han 
                    experimentado una evolución secular mediada por la 
                    relajación colisional, e.g., colapso del núcleo, segregación 
                    de masa y evaporación estelar 
                    \cite{chandrasekhar2005principles}. No obstante, el tiempo 
                    de cruce orbital sigue siendo muy inferior a 
                    $t_\text{relax}$, permitiendo que las ecuaciones de Jeans 
                    describan acertadamente el equilibrio instantáneo 
                    cuasi-estático sobre el cual operan los lentos procesos 
                    colisionales \cite{hamilton_2021_kinetic}.

                    \item \textbf{Cúmulos Abiertos:} Entornos escasamente 
                    poblados, $N \sim 10^2 - 10^4$, que habitan en los discos 
                    galácticos. Su $t_\text{relax}$ puede ser del orden de 
                    $10^7$ años \cite{clarke1999formation}. Representan el 
                    límite desafiante del análisis de Jeans. Las formulaciones 
                    no colisionales deben emplearse con cautela y acompañarse de 
                    simulaciones de N-cuerpos directas, ya que las colisiones 
                    proporcionan la semilla ruidosa inicial de interacciones 
                    rápidas y disrupción por el campo de marea galáctico 
                    \cite{bamber2025evolution}.
                \end{itemize}

            \subsubsection{Equilibrio Dinámico}
                La aplicación práctica de la teoría asume que el sistema ha 
                alcanzado un estado estacionario, lo que matemáticamente se 
                traduce en que la derivada parcial respecto al tiempo de 
                cualquier propiedad macroscópica es idénticamente nula, 
                $\partial / \partial t = 0$ \newline 
                \cite{binney_tremaine_2008}. Un sistema estelar naciente alcanza 
                este equilibrio a través de un mecanismo conocido como 
                relajación violenta, postulado por Donald Lynden-Bell, en el 
                cual las fluctuaciones macroscópicas y rápidas del potencial 
                gravitacional global termalizan la distribución estelar en unos 
                pocos tiempos de cruce. Para aplicar las ecuaciones de Jeans a 
                un cúmulo estelar, es indispensable que este no se encuentre en 
                un estado de fusión con otro sistema, ni sufriendo un desgarre 
                severo de marea reciente que distorsione gravemente su estado 
                estacionario subyacente \cite{mamon2010kinematic}.

            \subsubsection{Relación con el Potencial Gravitacional}
                La densidad generadora de gravedad y la densidad del trazador 
                cinemático no son intrínsecamente equivalentes 
                \cite{mamon2010kinematic}. Las ecuaciones de Jeans modelan un 
                "trazador", representado por una densidad espacial poblacional 
                $\nu(\mathbf{x})$, que se asume en equilibrio dentro del 
                potencial gravitacional $\Phi(\mathbf{x})$. Este potencial 
                obedece a la Ecuación de Poisson generalizada 
                \cite{binney_tremaine_2008}:

                $$\nabla^2 \Phi = 4\pi G \rho_\text{tot}$$

                donde $\rho_{tot}$ es la densidad de masa total. En cúmulos 
                estelares simples, $\rho_{tot}$ es esencialmente equivalente a 
                la masa de las estrellas, es decir, $\rho_{tot} = \Upsilon \nu$, 
                donde $\Upsilon$ es la razón masa-luminosidad 
                \cite{wolf_etal_2010}. Sin embargo, en galaxias esferoidales 
                enanas o cúmulos inmersos en densas concentraciones bariónicas y 
                materia oscura, el campo gravitacional total está dominado por 
                constituyentes invisibles ($\rho_{tot} \neq \nu$), haciendo que 
                la resolución de las ecuaciones de Jeans sea un problema de 
                inferencia paramétrica del perfil de masa oculta 
                \cite{walker2007velocity}.

        \subsection{Derivación de las Ecuaciones de Jeans y los Momentos de 
        Velocidad}
            Para transformar la CBE del espacio de fase al espacio de 
            configuración, se calcula la integral de la ecuación sobre todo el 
            dominio local de velocidades. Esto genera una serie infinita de 
            ecuaciones interconectadas conocida como la jerarquía de Jeans 
            \cite{ciotti_2021_jeans}. A nivel teórico, esta jerarquía nunca se 
            cierra por sí sola; el momento de orden $n$ siempre depende del 
            momento de orden $n+1$. Para hacerla tratable, la jerarquía se 
            trunca en el segundo momento introduciendo hipótesis físicas sobre 
            la asimetría de la varianza.

            Definimos la densidad espacial estelar (momento de orden cero de 
            $f$) como:

            $$\nu(\mathbf{x}, t) = \int f(\mathbf{x}, \mathbf{v}, t) d^3v$$

            La velocidad media o campo de velocidades sistemático local 
            (momento de primer orden) es: \cite{binney_tremaine_2008}

            $$\bar{v}_i(\mathbf{x}, t) = \frac{1}{\nu} \int v_i f(\mathbf{x}, 
            \mathbf{v}, t) d^3v$$

            \subsubsection{Ecuación de Momento de Orden Cero}
                Integrando la CBE término a término sobre el dominio de 
                velocidad $d^3v$:

                $$\int \frac{\partial f}{\partial t} d^3v + \int v_i 
                \frac{\partial f}{\partial x_i} d^3v - \int \frac{\partial 
                \Phi}{\partial x_i} \frac{\partial f}{\partial v_i} d^3v = 0$$

                Dado que la integración es sobre $\mathbf{v}$ y las coordenadas 
                $\mathbf{x}$ y $t$ son independientes, las derivadas espaciales 
                y temporales pueden salir de la integral. El tercer término se 
                anula aplicando el teorema de la divergencia sobre el espacio de 
                velocidades, bajo el pre-requisito físico de que $f \to 0$ con 
                la suficiente rapidez cuando $|\mathbf{v}| \to \infty$, el 
                sistema está gravitacionalmente ligado y no posee partículas de 
                energía infinita. El resultado es la ecuación de continuidad 
                \cite{binney_tremaine_2008}:

                $$\frac{\partial \nu}{\partial t} + \frac{\partial (\nu 
                \bar{v}_i)}{\partial x_i} = 0$$

                Esta expresión dicta la conservación del número de estrellas: 
                cualquier variación temporal en la densidad espacial estelar en 
                una región está forzosamente compensada por el flujo neto o 
                divergencia de la velocidad media en las fronteras de dicha 
                región.

            \subsubsection{Ecuación de Momento de Primer Orden}
                Para derivar la ecuación central que vincula la dispersión de 
                velocidades con el gradiente gravitacional, multiplicamos la CBE 
                por la componente de velocidad $v_j$ y repetimos la integración:

                $$\int v_j \frac{\partial f}{\partial t} d^3v + \int v_j v_i 
                \frac{\partial f}{\partial x_i} d^3v - \int v_j \frac{\partial 
                \Phi}{\partial x_i} \frac{\partial f}{\partial v_i} d^3v = 0$$

                La integral del último término, empleando integración por partes 
                y asumiendo nuevamente que $f$ se desvanece en el infinito, 
                produce un conmutador tensorial:

                $$- \int v_j \frac{\partial f}{\partial v_i} d^3v = \int f 
                \frac{\partial v_j}{\partial v_i} d^3v = \int f \delta_{ij} d^3v 
                = \nu \delta_{ij}$$

                Por ende, el término dependiente del potencial se contrae a $\nu 
                \frac{\partial \Phi}{\partial x_j}$. Seguidamente, descomponemos 
                el término cuadrático de velocidad $v_i v_j$ empleando la 
                velocidad media $\bar{v}_i$ y las velocidades aleatorias 
                intrínsecas locales. Definimos el tensor espacial de dispersión 
                de velocidades $\sigma_{ij}^2$, el cual mide la varianza 
                cinemática del campo:

                $$\sigma_{ij}^2 \equiv \overline{v_i v_j} - \bar{v}_i \bar{v}_j 
                = \frac{1}{\nu} \int (v_i - \bar{v}_i)(v_j - \bar{v}_j) f d^3v$$

                Sustituyendo $\overline{v_i v_j} = \sigma_{ij}^2 + \bar{v}_i 
                \bar{v}_j$ en las integrales, y utilizando la ecuación de 
                continuidad para cancelar componentes derivativos acoplados, 
                obtenemos la Ecuación de Jeans \cite{binney_tremaine_2008}:

                $$\nu \frac{\partial \bar{v}_j}{\partial t} + \nu \bar{v}_i 
                \frac{\partial \bar{v}_j}{\partial x_i} = -\nu \frac{\partial 
                \Phi}{\partial x_j} - \frac{\partial (\nu \sigma_{ij}^2)}{
                \partial x_i}$$

        \subsection{Interpretación Física y Analogía Hidrodinámica}
                El significado físico de la ecuación de Jeans adquiere profunda 
                elocuencia cuando se compara con la ecuación de movimiento de 
                Euler o Navier-Stokes, sin viscosidad cinemática que rige la 
                dinámica de los gases.

                El lado izquierdo de la ecuación de Jeans, $\nu (\frac{\partial 
                \bar{v}_j}{\partial t} + \bar{v}_i \frac{\partial \bar{v}_j}{
                \partial x_i})$, representa la aceleración convectiva de un 
                volumen infinitesimal del "fluido" de trazadores estelares a lo 
                largo de las trayectorias de su flujo sistemático medio 
                \cite{binney_tremaine_2008}.

                En el lado derecho de la ecuación residen las fuerzas 
                macroscópicas por unidad de volumen  El término $-\nu \nabla 
                \Phi$ es puramente la fuerza de atracción generada por el pozo 
                de gravedad. El aspecto crucial reside en el último término: 
                $-\frac{\partial (\nu \sigma_{ij}^2)}{\partial x_i}$. En un 
                fluido ordinario, el equivalente a este término es el gradiente 
                de presión isotrópica térmica $-\nabla P$ 
                \cite{binney_tremaine_2008}. En un sistema estelar no colisional 
                puro, la resistencia implosiva frente a la gravedad aplastante 
                no deviene de fuerzas repulsivas intermoleculares, sino 
                exclusivamente de la varianza estocástica geométrica de las 
                órbitas estelares. Por ello, al producto $\nu \sigma_{ij}^2$ se 
                le denomina el tensor de presión dinámica anisótropa 
                \cite{mamon2010kinematic}.

                Sin embargo, a diferencia de un gas molecular termalizado por 
                frecuentes choques donde la dispersión es estadísticamente 
                idéntica en cualquier eje, las estrellas en un cúmulo o halo 
                retienen memorias orbitales dictadas por sus constantes de 
                movimiento, de forma tal que $\sigma_{ij}^2$ no tiene por qué 
                ser diagonal ni isotrópico. La presión estelar puede tener 
                dependencias preferenciales o de cizallamiento a través de la 
                presencia de términos no diagonales \cite{binney_tremaine_2008}. 
                Si tomamos el primer momento espacial general (multiplicar toda 
                la ecuación por las coordenadas $x_k$ e integrar sobre el 
                volumen general del cúmulo), se recupera el Teorema del Virial 
                en forma tensorial, el cual es el principio gobernante 
                subyacente para las inferencias iniciales globales sobre masa en 
                dinámica galáctica \cite{binney_tremaine_2008}.

        \subsection{Formulación Esférica de la Ecuación de Jeans}
                El tratamiento del tensor de dispersión de velocidades 
                $\sigma_{ij}^2$ requiere adoptar sistemas de coordenadas 
                alineados a las simetrías inherentes del modelo astronómico, 
                reduciendo las seis variables teóricas independientes a un 
                número conmensurable con las observables. A continuación, se 
                detalla la formulación esférica de la ecuación de Jeans.

                \subsubsection{Ecuación Esférica de Jeans}
                    El marco de aplicación por excelencia para los cúmulos 
                    estelares globulares y para cúmulos abiertos altamente 
                    masivos es la aproximación esférica. Evaluada en coordenadas 
                    polares esféricas $(r, \theta, \phi)$, la simetría fuerza a 
                    que el potencial y la densidad sean únicamente funciones del 
                    radio $r$, anulando la dependencia azimutal e inclinación 
                    \cite{mamon2010kinematic}.

                    Por argumentos de simetría fundamental invariante bajo 
                    rotación, el tensor de esfuerzos cruzados se anula 
                    ($\sigma_{r\theta}^2 = \sigma_{r\phi}^2 = \sigma_{
                    \theta\phi}^2 = 0$) \cite{courteau2014galaxy}, asumiendo que 
                    los ejes principales del elipsoide de velocidades están 
                    perfectamente alineados con el sistema coordenado. Asumiendo 
                    estacionalidad ($\partial/\partial t = 0$) y ausencia de 
                    corrientes netas ($\bar{v} = 0$), el complejo sistema 
                    matricial se colapsa en una única ecuación diferencial 
                    ordinaria hidrostática que vincula la presión radial con la 
                    componente gravitatoria \cite{binney_tremaine_2008}: 

                    $$\frac{d(\nu \sigma_r^2)}{dr} + \frac{2\nu}{r}\left(
                    \sigma_r^2 - \frac{\sigma_\theta^2 + \sigma_\phi^2}{2} 
                    \right) = -\nu \frac{d\Phi}{dr}$$

                    Dada la conservación de momento angular isótropa en una 
                    morfología esférica media, las dispersiones tangenciales 
                    angulares acaban siendo estocásticamente indiferenciables: 
                    $\sigma_\theta^2 = \sigma_\phi^2 \equiv \sigma_t^2$. 
                    Reestructurando el gradiente gravitacional mediante la masa 
                    contenida newtoniana $M(r)$, la expresión simplificada es 
                    \cite{mamon2010kinematic}:

                    $$\frac{1}{\nu}\frac{d(\nu \sigma_r^2)}{dr} + \frac{2}{r}(
                    \sigma_r^2 - \sigma_t^2) = - \frac{GM(r)}{r^2}$$

                    Asumiendo un estado de isotropía total en la dispersión de 
                    velocidades, la dispersión es idéntica en todas las 
                    direcciones del sistema, es decir, $\sigma_r^2 = \sigma_t^2 
                    \equiv \sigma^2$. Con esta consideranción, el término de 
                    anisotropía cruzada se anula, reduciendo la ecuación 
                    anterior a la forma:

                    $$\frac{1}{\nu} \frac{d(\nu \sigma^2)}{dr} = 
                    -\frac{GM(r)}{r^2}$$

                    Esta ecuación diferencial ordinaria permite una solución 
                    mediante integración directa desde un radio interior de 
                    interés $r$ hasta el infinito, dando lugar a la formulación 
                    integral de la ecuación de Jeans 
                    \cite{binney_tremaine_2008}:

                    \begin{equation}
                        \sigma^2(r) = \frac{1}{\nu(r)} \int_r^\infty \nu(r') 
                        \frac{GM(r')}{r'^2} dr'
                        \label{eq:jeans_integral}
                    \end{equation}
       
            \subsection{Dispersión de Velocidades y Anisotropía}
                El término divergente en la ecuación esférica de Jeans ha 
                introducido formalmente las componentes vectoriales de la 
                varianza de la velocidad orbital: la dispersión radial 
                $\sigma_r$ y la dispersión tangencial bidimensional $\sigma_t$ 
                \cite{wolf_etal_2010}. La relación diferencial subraya un pilar 
                físico de extrema trascendencia: a masa dinámica no puede 
                inducirse exclusivamente a partir de las dispersiones 
                individuales sin entender la correlación geométrica conjunta que 
                adoptan las estrellas orbitantes.

                Esta conjunción orbital se parametriza universalmente a través 
                del parámetro de anisotropía de velocidad de Binney, denotado 
                como $\beta(r)$ \cite{binney_tremaine_2008}:

                $$\beta(r) \equiv 1 - \frac{\sigma_\theta^2 + \sigma_\phi^2}{2
                \sigma_r^2} = 1 - \frac{\sigma_t^2}{\sigma_r^2}$$

                El parámetro $\beta(r)$ es un indicador astrofísico adimensional 
                que clasifica el ensamble topológico de las órbitas dominantes 
                en la capa esférica del sistema estelar a una profundidad $r$ 
                \cite{binney_tremaine_2008}:

                \begin{itemize}
                    \item $\beta = 0$: El sistema es un rotador isotrópico de 
                    fase espacial. La presión gravitacional expansiva es 
                    ejercida equivalentemente en todas las direcciones 
                    tridimensionales $\sigma_r = \sigma_t$.

                    \item $0 < \beta \leq 1$: Las estrellas ocupan un estado 
                    anisótropo con sesgo radial. Las partículas se sumergen 
                    gravitacionalmente hacia el pozo profundo de núcleo y 
                    vuelven a alejarse en cuerdas elongadas. En el límite ideal 
                    $\beta = 1$, las órbitas son puramente puras agujas radiales 
                    y $\sigma_t = 0$.

                    \item $\beta < 0$: El clúster adolece de un sesgo 
                    tangencial. Las órbitas estelares son predominantemente 
                    circulares, análogas a discos difusos. A medida que $\beta 
                    \to -\infty$, $\sigma_r \to 0$.
                \end{itemize}

                La ecuación de Jeans esférica definitiva en términos de este 
                coeficiente se consolida finalmente como:

                \begin{equation}
                    \frac{d}{dr}(\nu \sigma_r^2) + \frac{2\beta(r)}{r} \nu 
                    \sigma_r^2 = - \nu \frac{GM(r)}{r^2}
                    \label{eq:jeans_esferica}
                \end{equation}

                Dado que el parámetro de anisotropía requiere conocimiento 
                intrínseco de la función de distribución del espacio de fase 
                inobservable $f(E, J)$, donde $E$ es la energía orbital y $J$ es 
                el momento angular, investigaciones adoptan modelos heurísticos 
                pre-configurados para $\beta(r)$ \cite{mamon2010kinematic}.
